\chapter{Strings e F-Strings}

\section{Literais de string}

Strings são delimitadas por aspas duplas:

\begin{lstlisting}
let s = "Olá, mundo"
let caminho = "/home/usuario/dados.csv"
let vazia = ""
\end{lstlisting}

\subsection*{Sequências de escape}

\begin{center}
\begin{tabular}{cl}
\toprule
\textbf{Sequência} & \textbf{Significado} \\
\midrule
\texttt{\textbackslash n}  & Nova linha \\
\texttt{\textbackslash t}  & Tabulação \\
\texttt{\textbackslash r}  & Retorno de carro \\
\texttt{\textbackslash"}   & Aspas duplas literais \\
\texttt{\textbackslash\textbackslash} & Barra invertida literal \\
\bottomrule
\end{tabular}
\end{center}

\section{Operações com strings}

\subsection*{Concatenação}

\begin{lstlisting}
let nome = "Hayashi"
let versao = "0.2"
let titulo = nome + " " + versao    // "Hayashi 0.2"
\end{lstlisting}

\subsection*{Comprimento}

\begin{lstlisting}
let s = "econometria"
display len(s)    // 11
\end{lstlisting}

\subsection*{Conversões}

\begin{lstlisting}
let n = str(42)        // "42"
let f = str(3.14)      // "3.14"
let b = str(true)      // "true"
\end{lstlisting}

\subsection*{Comparação}

Strings são comparadas lexicograficamente com \lstinline{==}, \lstinline{!=},
\lstinline{<}, \lstinline{>}, \lstinline{<=}, \lstinline{>=}.

\section{F-Strings}

F-strings permitem interpolar expressões arbitrárias dentro de uma string. O
prefixo \lstinline{f} indica interpolação; expressões entre \lstinline|{ }| são
avaliadas e convertidas para string:

\begin{lstlisting}[caption={F-strings}]
let municipio = "Pelotas"
let pop = 328_000

display f"{municipio} tem {pop} habitantes"
// Pelotas tem 328000 habitantes

display f"Média: {(10 + 20 + 30) / 3.0:.2f}"
// Média: 20.00
\end{lstlisting}

\subsection*{Expressões em f-strings}

Qualquer expressão pode ser interpolada — variáveis, operações aritméticas,
chamadas de função:

\begin{lstlisting}
let n = 5
display f"{n}^2 = {n^2}"
// 5^2 = 25

display f"raiz de {n} = {sqrt(n):.4f}"
// raiz de 5 = 2.2361
\end{lstlisting}

\subsection*{F-strings em \texttt{display}}

\lstinline{display} aceita opções \lstinline{end=} e \lstinline{sep=} para controlar
o terminador e o separador entre múltiplos valores:

\begin{lstlisting}
display "a", "b", "c" sep=", "   // a, b, c
display "linha" end=""            // linha (sem quebra)
\end{lstlisting}

\section{Multiline strings}

\hayashi{} não tem literal de string multiline. Use concatenação explícita:

\begin{lstlisting}
let sql = "SELECT * FROM dados " +
          "WHERE ano = 2026 " +
          "ORDER BY uf"
\end{lstlisting}
